\documentclass [12pt] {article}
\usepackage{amsmath, amssymb}

\title{Inductive proof of Nicomachus's Theorem}
\author{Heejun Yoon}
\date{\today}

\begin{document}

\maketitle

\begin{abstract}
    This paper shows proof by induction of Nicomachus's Theorem in high school level mathematics.
\end{abstract}

\tableofcontents
\newpage 

\section{Introduction}
While Nicomachus's theorem is a classical result in number theory, it may be applied for efficient calculation for computers.

\section{Nicomachus's Theorem}
Nicomachus's Theorem shows that the addition of $1^3, 2^3 \dots n^3$ is identical to the addition of $1$ to $n$ squared. This is written as: 
\begin{equation}
    \sum_{i=1}^{n} i^3 = \left(\sum_{i=1}^{n}i\right)^2
\end{equation}

\section{Induction}
The statement P(n) is defined as for $n \in \mathbb{Z}^+$,
\begin{equation}
    \sum_{i=1}^{n} i^3 = \left(\sum_{i=1}^{n}i\right)^2
\end{equation}
holds.

\subsection{Base case}
when $n=1$, 
\begin{align*}
    \text{LHS} &=\sum_{i=1}^{1}i^3\\
    &=1^3\\
    &=1
\end{align*}
\begin{align*}
    \text{RHS} &=\left(\sum_{i=1}^{1}i\right)^2\\
    &=(1)^2\\
    &=1
\end{align*}
Therefore, $\text{LHS}=\text{RHS}$, base case $n=1$ holds.

\subsection{Inductive Hypothesis}
Assume the statement P(k), for $k \in \mathbb{Z}^+$,
\begin{equation}
    \sum_{i=1}^{k} i^3 = \left(\sum_{i=1}^{k}i\right)^2
\end{equation}
holds.

\subsection{Proof by induction}
Prove that the statement P(k+1) holds.
\begin{equation}
    \sum_{i=1}^{k+1} i^3 = \left(\sum_{i=1}^{k+1}i\right)^2
\end{equation}
\begin{align*}
    \text{RHS} &= \left(\sum_{i=1}^{k+1}i\right)^2 \\
    &= \left\{\sum_{i=1}^{k}i + (k+1)\right\}^2 \\
    &= \left(\sum_{i=1}^ki\right)^2 + 2(k+1) \times \sum_{i=1}^ki + (k+1)^2 \\
    &= \sum_{i=1}^{k}i^3 + (k+1)\left(k+1+2 \times \sum_{i=1}^{k}i\right) && \text{by 3.2} \\
    &= \sum_{i=1}^{k}i^3 + (k+1)\left\{k+1+2 \times \left(\frac{k(k+1)}{2}\right)\right\} && \text{by sum of arithmetic sequence} \\
    &= \sum_{i=1}^{k}i^3 + (k+1)(k+1+k^2+k) \\
    &= \sum_{i=1}^{k}i^3 + (k+1)(k^2+2k+1) \\
    &= \sum_{i=1}^{k}i^3 + (k+1)(k+1)^2 \\
    &= \sum_{i=1}^{k}i^3 + (k+1)^3 \\
    &= \sum_{i=1}^{k+1}i^3 \\[0.5em]
    \therefore \text{LHS} &= \text{RHS} &&\hfill \square
\end{align*}

\section{Conclusion}
Hence, according to the Principle of Mathematical Induction, since the base case and the inductive step are both proven, the statement holds for all \(n \in \mathbb{Z}^+\).

\end{document}
